\documentclass[a4paper,10pt]{report}

\def\packagepath{../../preambule}   % path du package princial
\usepackage{\packagepath/preambule}    % utilisation du fichier de configuration

\def\level{Première }              % Classe
\def\course{NSI}              % Matière
\def\eval{DS04}

\def\visibleornot{visible}    % visible or invisible
\def\documentpath{./Sources_Latex}
\def\date{Jeudi 06 Octobre 2025}
\renewcommand{\arraystretch}{1}  % Ecart dans les tableau

\def\ptsexoA{2}
\def\ptsexoB{1}
\def\ptsexoC{1}
\def\ptsexoD{1}
\def\ptsexoE{1}
\def\ptsexoF{1}
\def\ptsexoG{2}
\def\ptsexoH{1}
\def\ptsexoI{2}
\def\ptsexoJ{2}
\def\ptsexoK{1}

\def\ptstotal{\ptsexoA+\ptsexoB}

\usepackage{hyperref}

\usetikzlibrary{shapes, arrows, positioning}
\usetikzlibrary{shapes.geometric, arrows, positioning, decorations.pathreplacing}

\tikzstyle{etat} = [draw, rounded corners=15pt, minimum width=2.5cm, minimum height=1.2cm, text centered, font=\bfseries]
\tikzstyle{nouveau} = [etat, fill=blue!40]
\tikzstyle{pret} = [etat, fill=green!60]
\tikzstyle{elu} = [etat, fill=yellow!60]
\tikzstyle{bloque} = [etat, fill=red!60]
\tikzstyle{termine} = [etat, fill=white]
\tikzstyle{fleche} = [thick, ->, >=stealth]

\usepackage{float} % Permet d'utiliser [H]
\usepackage[T1]{fontenc}
\begin{document}

%%%%%%%%%%%%%%%%%%%%%%%%%%%%%%%%%%%%%%%%%%%%%%%%%%ù

%\PageGardeSujetBac[Matiere = Numérique et Sciences Informatiques,
%                  Session = 2025,
%                  AffJour=false,
%                  Duree = {3 heures 30},
%                  DernierePage = \pageref{LastPage},
%                  ModeExamen=false
%                  ]
\renewcommand{\labelitemi}{\textbullet} %pour éviter les tirets dans les "itemize" qui apportent confusion avec le signe moins.
%% Début page de garde style BAC

   %%%%%%%%%%%%%%%%%%%%%%%%%%%%%%%%%%%%%%%%%%%%%%%%%%%%%%%%
%\PageGardeSujetBac[clés]

\pagestyle{DS_FP}          % Feuille de style fancy
\NomPrenomNote{}


\exods{\ptsexoA}

Sans aucune justification, donner le nombre d'états possibles codables avec:

\begin{enumerate}[1)]
    \item 1 bit \dotfill{}
    \item 3 bits \dotfill{}
    \item 8 bits \dotfill{}
    \item n bits \dotfill{}
\end{enumerate}

\medskip


\begin{minipage}{0.49\linewidth}
 \exods{\ptsexoB}

   Convertir le nombre binaire suivant en décimal. Vous justifierez vos réponses.

\begin{answer}{1}{5} 
    \begin{center}
        \verb|1010 0011|
    \end{center}   
\end{answer}
\end{minipage}\hfill
\begin{minipage}{0.49\linewidth}
     \exods{\ptsexoC}

    Convertir le nombre décimal suivant en binaire. Vous justifierez vos réponses.
\begin{answer}{1}{5} 
    \begin{center}
        \verb|159|
    \end{center}   
\end{answer}
\end{minipage}

\medskip


\begin{minipage}{0.49\linewidth}
 \exods{\ptsexoD}

   Convertir le nombre hexadécimal suivant en décimal. Vous justifierez vos réponses.

\begin{answer}{1}{5} 
    \begin{center}
        \verb|AF|
    \end{center}   
\end{answer}
\end{minipage}\hfill
\begin{minipage}{0.49\linewidth}
     \exods{\ptsexoE}

Convertir le nombre binaire suivant en hexadécimal. Vous justifierez vos réponses.
\begin{answer}{1}{5} 
    \begin{center}
        \verb|1110 0110|
    \end{center}   
\end{answer}
\end{minipage}

\medskip

\exods{\ptsexoF}

Effectuer l'addition binaire suivante:



    \begin{answer}{1}{4} \begin{center}
        \verb|10110111+1101110|
\end{center}  \end{answer}

\pagebreak

\exods{\ptsexoG}

Expliquer simplement comment déterminer si un nombre binaire est pair ou impair. Justifier.

    \begin{answer}{1}{3}    \end{answer}




\medskip

\exods{\ptsexoH}

Ecrire une fonction \verb|est_pair(n)| qui prend en paramètre un nombre écrit en binaire (au format chaine de caractères) et renvoie \verb|True| si le nombre est pair et \verb|False| sinon.

    \begin{answer}{1}{7} 
\begin{verbatim}
def est_pair(n):
    ''' Renvoie True si n est pair et False sinon
    n chaîne de caractères représentant un nombre binaire
    sortie: bool '''









assert est_pair('1010') == True    # 10 en décimal
assert est_pair('111') == False    # 7 en décimal
    \end{verbatim}

    \end{answer}

\medskip

\exods{\ptsexoI}

Ecrire une fonction \verb|nb_bits(n)| qui prend en paramètre un entier naturel \verb|n| et renvoie le nombre de bits nécessaires pour l'écrire en binaire. Par exemple, \verb|nb_bits(5)| renvoie \verb|3| car $5 = 101_2$.

On pourra par exemple compter le nombre de divisions entières par 2 nécessaires pour obtenir 0.
    \begin{answer}{1}{9.2} 
\begin{verbatim}
def nb_bits(n):
    ''' Renvoie le nombre de bits nécessaires pour écrire n en binaire 
    n entier naturel (int)
    sortie: int '''














assert nb_bits(5) == 3    # 5 = 101_2
assert nb_bits(10) == 4   # 10 = 1010_2 
assert nb_bits(16) == 5   # 16 = 10000_2
    \end{verbatim}

    \end{answer}

    \pagebreak

    \exods{\ptsexoJ}

Ecrire une fonction \verb|nb_uns(val_bin)| qui prend en paramètre une chaine de caractères \verb|val_bin| représentant un nombre binaire et qui renvoie le nombre de bits à 1 dans cette chaine.

    \begin{answer}{1}{9}
\begin{verbatim}
def nb_uns(val_bin):
    ''' Renvoie le nombre de bits à 1 dans val_bin
    val_bin chaine de caractères représentant un nombre binaire
    













assert nb_uns('101010') == 3
assert nb_uns('1111') == 4
assert nb_uns('0000') == 0
    \end{verbatim}

    \end{answer}

    \medskip

     \exods{\ptsexoK}

Ecrire une fonction \verb|multiplication(val_bin, puissance)| qui va multiplier un nombre binaire \verb|val_bin| (sous forme de chaine de caractères) par $2^{puissance}$ où \verb|puissance| est un entier naturel.
    \begin{answer}{1}{10}
\begin{verbatim}
def multiplication(val_bin, puissance):
    ''' Multiplie le nombre binaire val_bin par 2^puissance
    val_bin chaine de caractères représentant un nombre binaire
    puissance entier naturel (int)
    sortie: str '''














assert multiplication('101', 3) == '101000'    # 5 * 2^3 = 40
assert multiplication('111', 2) == '11100'     # 7 * 2^2 = 28
assert multiplication('1001', 0) == '1001'     # 9 * 2^0 = 9
    \end{verbatim}

    \end{answer}

    \pagebreak
    \thispagestyle{DS_LP}

\exods{0}

    On dispose de deux chaines binaires \verb|A| et \verb|B| de même longueur (8 bits).

    On veut produire une chaine binaire représentant leur somme.

    Ecrire une fonction \verb|addition_binaire(A, B)| qui prend en paramètres les deux chaines binaires \verb|A| et \verb|B| et renvoie la chaine binaire représentant leur somme. On ne pourra pas convertir les valeurs en décimal avant de faire l'addition.

\begin{answer}{1}{11} 
\begin{verbatim}
def addition_binaire(A, B):
    ''' Renvoie la chaine binaire représentant la somme de A et B
    A et B chaines de caractères représentant des nombres binaires de même longueur (8 bits)
    sortie: str '''















assert addition_binaire('00011010', '00000101') == '00011111'    # 26 + 5 = 31
assert addition_binaire('11111111', '00000001') == '100000000'    # 255 + 1 = 256
assert addition_binaire('10101010', '01010101') == '11111111'    # 170 + 85 = 255
    \end{verbatim}

    \end{answer}
\end{document}

    \medskip

    

On exécute le script python ci-dessous, que vaut le couple \verb|(a, b)| à la fin de l'exécution?


\begin{CodePythontexAlt}[TaillePolice=\large,Largeur=1\linewidth,Centre,PremLigne=1]{}
def carre(a):
    a = a * a
    return a

a = 2
b = carre(a)
\end{CodePythontexAlt}


    \begin{answer}{1}{1}    \end{answer}

\medskip

\exods{\ptsexoB}

La fonction \verb|ajoute(n, p)| calcule la somme de tous les entiers entre \verb|n| et \verb|p| inclus. On suppose \verb|n < p|.

Compléter cette fonction

\begin{CodePythontexAlt}[TaillePolice=\large,Largeur=1\linewidth,Centre,PremLigne=1]{}
def ajoute(n, p):
    ''' Ajoute tous les entiers entre n et p inclus
    n et p entiers (int) et n < p
    sortie: entier '''
    somme = ...... 
    for i in range( ...... , ......):
        somme = ............
    return ............ 


assert ajoute(2, 4) == 9    # 2 + 3 + 4 = 9
    \end{CodePythontexAlt}


    \medskip

\exods{\ptsexoC}

Ecrire une fonction qui renvoie \verb|True| si \verb|a| est strictement supérieur à \verb|b| et \verb|False| sinon.

\begin{CodePythontexAlt}[TaillePolice=\large,Largeur=1\linewidth,Centre,PremLigne=1]{}
def compare(a, b):
    ''' Renvoie True si a strictement supérieur à b et False sinon
    sortie: bool '''







assert compare(1, 5) == False
assert compare(4, 2) == True
    \end{CodePythontexAlt}

 \pagebreak
    \pagestyle{DS_LP}

\exods{\ptsexoD}

Ecrire une fonction \verb|est_present(lettre_cherchee, mot)| qui renvoie \verb|True| si \verb|lettre_cherchee| est dans \verb|mot|, \verb|False| sinon

\begin{CodePythontexAlt}[TaillePolice=\large,Largeur=1\linewidth,Centre,PremLigne=1]{}
def est_present(lettre_cherchee, mot):
    ''' Renvoie True si lettre_cerchee est dans mot et False sinon / sortie: bool '''





assert est_present('a', 'bonjour') == False
assert est_present('o', 'bonjour') == True
assert est_present('n', 'bonjour') == True

    \end{CodePythontexAlt}

    \medskip

    \exods{\ptsexoE}


Ecrire une fonction \verb|compte_lettre_mot(lettre_cherchee, mot)| qui compte le nombre de \verb|lettre_cherchee| dans \verb|mot|

\begin{CodePythontexAlt}[TaillePolice=\large,Largeur=1\linewidth,Centre,PremLigne=1]{}
def compte_lettre_mot(lettre_cherchee, mot):
    ''' Compte le nombre de lettre_cherchee dans mot / sortie: int '''








assert compte_lettre_mot('a', 'bonjour') == 0
assert compte_lettre_mot('o', 'bonjour') == 2
assert compte_lettre_mot('n', 'bonjour') == 1
    \end{CodePythontexAlt}

    \medskip

    \exods{\ptsexoF}


Ecrire une fonction \verb|supprime_voyelles(mot)| qui renvoie le \verb|mot| passé en paramètre auquel on a supprimé les voyelles.

\begin{CodePythontexAlt}[TaillePolice=\large,Largeur=1\linewidth,Centre,PremLigne=1]{}

VOYELLES = 'aeiouy'

def supprime_voyelles(mot):
    ''' Renvoie le mot entré en paramètre sans les voyelles / sortie: str '''








assert supprime_voyelles('bonjour') == 'bnjr'
assert supprime_voyelles('html') == 'html'
assert supprime_voyelles('eau') == ''
    \end{CodePythontexAlt}


\end{document}


Indiquer si les noms de variables suivants sont pertinents ou non. Justifier votre réponse.

\medskip

\renewcommand{\arraystretch}{3.5}
\begin{tabular}{|p{3cm}|p{3cm}|p{3cm}|p{7cm}|}
        \hline
        \textbf{Nom de variable} & \textbf{Pertinent} & \textbf{Non pertinent} &\textbf{Justification} \\ \hline
    \verb|ma variable| &&& \\ \hline
    \verb|ma_variable77|&&& \\ \hline
    \verb|77ma_variable|&&& \\ \hline
    \verb|ma_variable@77|&&& \\ \hline

    \end{tabular}



\bigskip

\exods{\ptsexoB}

    \renewcommand{\arraystretch}{1}

Indiquer l'état des variables après l'exécution de chaque ligne d'instruction.

\medskip

\begin{minipage}{0.29\linewidth}

\begin{CodePythontexAlt}[TaillePolice=\large,Largeur=1\linewidth,Centre,PremLigne=1]{}
a = 3
b = 5
c = a + b
a = b - 2
b = c + 1
d = a + b + c
\end{CodePythontexAlt}

\end{minipage}\hfill
\begin{minipage}{0.66\linewidth}
    
    \medskip

    \begin{tabular}{|c|c|c|c|c|}
        \hline
        Ligne& Valeur de a & Valeur de b &Valeur de c & Valeur de d\\ \hline
        1&3&&&\\\hline
        2&&&&\\\hline
        3&&&&\\\hline
        4&&&&\\\hline
        5&&&&\\\hline
        6&&&&\\\hline


    \end{tabular}
\end{minipage}

\bigskip

\exods{\ptsexoC}

Indiquer le type et la valeur de la variable \verb|a| après l'exécution de chacune des lignes. Si la ligne génère une erreur, indiquer "erreur".


\begin{minipage}{0.49\linewidth}

\begin{CodePythontexAlt}[TaillePolice=\large,Largeur=1\linewidth,Centre,PremLigne=1]{}
a = 3
a = a + 0.1
a = "Bonjour"
a = a + " tout le monde"
a = True
a = a + 1
a = input('Entrer un nombre : ')
a = '5'
a = a + '2'
a = int(a)
a = float(a)

\end{CodePythontexAlt}

\end{minipage}\hfill
\begin{minipage}{0.46\linewidth}
    
    \medskip

    \renewcommand{\arraystretch}{1.2}
    \begin{tabular}{|c|p{2cm}|p{5cm}|}
        \hline
        Ligne& Type de a & Valeur de a\\ \hline
        1& \verb|int|&3\\\hline
        2&&\\\hline
        3&&\\\hline
        4&&\\\hline
        5&&\\\hline
        6&&\\\hline
        7&&\\\hline
        8&&\\\hline
        9&&\\\hline
        10&&\\\hline
        11&&\\\hline


    \end{tabular}
\end{minipage}


\pagebreak
\pagestyle{DS_LP}
\exods{\ptsexoD}

On considère le code suivant:

\begin{CodePythontexAlt}[TaillePolice=\large,Largeur=1\linewidth,Centre,PremLigne=1]{}
mot = 'Numerique'
\end{CodePythontexAlt}


Indiquer ce qui est affiché par les instructions suivantes. Si cela génère une erreur, indiquer "erreur".

\renewcommand{\arraystretch}{1.2}
    \begin{tabular}{|p{3cm}|p{4cm}|}
        \hline
         Instruction & Affichage\\ \hline
         \verb|mot[0]|&\\\hline
        \verb|mot[2]| &\\\hline
        \verb|mot[4]| &\\\hline
        \verb|mot[8]|&\\\hline
        \verb|len(mot)|&\\\hline
        \verb|mot + 'nsi'| &\\\hline
        
        
    \end{tabular}\hfill
\renewcommand{\arraystretch}{1.2}
    \begin{tabular}{|p{3cm}|p{4cm}|}
        \hline
         Instruction & Affichage\\ \hline
        \verb|'nsi' + mot|&\\\hline
        \verb|'u' in mot|&\\\hline
        \verb|'n' in mot|&\\\hline
        \verb|int(mot)|&\\\hline
        \verb|mot * 2| &\\\hline
        \verb|mot[3] * 2| &\\\hline
        
        


    \end{tabular}

\bigskip

\exods{\ptsexoE}

On considère le code suivant:

\begin{CodePythontexAlt}[TaillePolice=\large,Largeur=1\linewidth,Centre,PremLigne=1]{}
a = True
b = False
\end{CodePythontexAlt}


Indiquer ce qui est affiché par les instructions suivantes. Si cela génère une erreur, indiquer "erreur".

\renewcommand{\arraystretch}{1.2}
    \begin{tabular}{|p{4cm}|p{3cm}|}
        \hline
         Instruction & Affichage\\ \hline
         \verb|a and b|&\\\hline
        \verb|a or b| &\\\hline
        \verb|not a| &\\\hline
        \verb|a and (not b)|&\\\hline
        
        
    \end{tabular}\hfill
\renewcommand{\arraystretch}{1.2}
    \begin{tabular}{|p{4cm}|p{3cm}|}
        \hline
         Instruction & Affichage\\ \hline
        \verb|a or (not b)| &\\\hline
        \verb|a and (a or b)|&\\\hline
        \verb|b or (a and b)|&\\\hline
        \verb|a and (not(a and b))|&\\\hline
    \end{tabular}
\bigskip


\exods{\ptsexoF}

On considère le code suivant:

\begin{CodePythontexAlt}[TaillePolice=\large,Largeur=1\linewidth,Centre,PremLigne=1]{}
a = 3
b = 4
\end{CodePythontexAlt}


Indiquer ce qui est affiché par les instructions suivantes. Si cela génère une erreur, indiquer "erreur".

\renewcommand{\arraystretch}{1.2}
    \begin{tabular}{|p{4cm}|p{3cm}|}
        \hline
         Instruction & Affichage\\ \hline
         \verb|a == 3|&\\\hline
        \verb|a != 3| &\\\hline
        \verb|a > 3| &\\\hline
        
        
        
    \end{tabular}\hfill
\renewcommand{\arraystretch}{1.2}
    \begin{tabular}{|c|p{4cm}|p{3cm}|}
        \hline
         Instruction & Affichage\\ \hline
        \verb|type(a) == str|&\\\hline
        \verb|a == 4 and a/0 == 1|&\\\hline
        \verb|a == 3 or a/0 == 1| &\\\hline
         \end{tabular}
\bigskip


\exods{\ptsexoG}

Proposer un code Python permettant:

\begin{enumerate}
    \item de demander à un utilisateur son année de naissance.
    \item de lui demander son prénom
    \item de lui afficher le message suivant: "Bonjour \texttt{<prenom>}, tu as \texttt{<age>} ans." en remplaçant \verb|<age>| et \verb|<prenom>| par l'âge et le prénom de l'utilisateur. Par exemple:  \quad \verb|Bonjour Thomas, tu as 17 ans.|
    \end{enumerate}
    
    \begin{answer}{1}{2.6}    \end{answer}

\end{document}








